% GENERATED FILE. Edit canonical lessons, then run npm run generate:books.
% canonical-source-hash: fnv1a64:9335a09b4930776f
% canonical-lessons: TE-C297-mondi, TE-C297-pisinari, TE-C297-curukaina, TE-C297-allari, TE-C297-nijayiti

\chapter{Stubborn, Stingy, Active, Naughty, Honesty}
\label{ch:te-stubborn-stingy-active-naughty-honesty}
\hlchaptermodality{297}

\begin{chapteropening}
\textbf{By the end of this chapter:} \emph{I can say stubborn, stingy, active, naughty and honesty.}

Complete the last lesson of chapter 297: I can say stubborn, stingy, active, naughty and honesty.
\end{chapteropening}

\section[\texorpdfstring{\te{మొండి} (moṇḍi)}{మొండి (moṇḍi)}]{\te{మొండి} (moṇḍi) — stubborn}

\label{lesson:TE-C297-mondi}



\begin{quote}
Before the new one: say the Telugu for useful, then the Telugu for local.
\end{quote}

\subsection*{You'll want to know: \te{మొండి}}
\textbf{\te{మొండి}} — \emph{moṇḍi} — \textquotedblleft{}stubborn\textquotedblright{}.

\begin{grammarlens}[title={What you've built}]
One more word for this chapter.
\end{grammarlens}

\subsection*{Your turn}
\begin{itemize}
\raggedright
  \item \emph{Say it:} \emph{moṇḍi}
  \item \emph{Say it:} \emph{moṇḍi}, once more
  \item \emph{Say it:} say \emph{moṇḍi}
  \item \emph{Recall:} say the Telugu for useful, then the Telugu for local, then say \emph{moṇḍi} again
\end{itemize}

\begin{tcolorbox}[breakable,colback=teal!4,colframe=teal!35!black,title={Before you move on}]
What does \te{మొండి} mean? (\textquotedblleft{}Stubborn\textquotedblright{}.) Say it once more.
\end{tcolorbox}
\section[\texorpdfstring{\te{పిసినారి} (pisināri)}{పిసినారి (pisināri)}]{\te{పిసినారి} (pisināri) — stingy; a miser}

\label{lesson:TE-C297-pisinari}



\begin{quote}
Before the new one: say the Telugu for local, then the Telugu for stubborn.
\end{quote}

\subsection*{You'll want to know: \te{పిసినారి}}
\textbf{\te{పిసినారి}} — \emph{pisināri} — \textquotedblleft{}stingy; a miser\textquotedblright{}.

\begin{grammarlens}[title={What you've built}]
One more word for this chapter.
\end{grammarlens}

\subsection*{Your turn}
\begin{itemize}
\raggedright
  \item \emph{Say it:} \emph{pisināri}
  \item \emph{Say it:} \emph{pisināri}, once more
  \item \emph{Say it:} say \emph{pisināri}
  \item \emph{Recall:} say the Telugu for local, then the Telugu for stubborn, then say \emph{pisināri} again
\end{itemize}

\begin{tcolorbox}[breakable,colback=teal!4,colframe=teal!35!black,title={Before you move on}]
What does \te{పిసినారి} mean? (\textquotedblleft{}Stingy; a miser\textquotedblright{}.) Say it once more.
\end{tcolorbox}
\section[\texorpdfstring{\te{చురుకైన} (curukaina)}{చురుకైన (curukaina)}]{\te{చురుకైన} (curukaina) — active, alert}

\label{lesson:TE-C297-curukaina}



\begin{quote}
Before the new one: say the Telugu for stubborn, then the Telugu for stingy.
\end{quote}

\subsection*{You'll want to know: \te{చురుకైన}}
\textbf{\te{చురుకైన}} — \emph{curukaina} — \textquotedblleft{}active, alert\textquotedblright{}.

\begin{grammarlens}[title={What you've built}]
One more word for this chapter.
\end{grammarlens}

\subsection*{Your turn}
\begin{itemize}
\raggedright
  \item \emph{Say it:} \emph{curukaina}
  \item \emph{Say it:} \emph{curukaina}, once more
  \item \emph{Say it:} say \emph{curukaina}
  \item \emph{Recall:} say the Telugu for stubborn, then the Telugu for stingy, then say \emph{curukaina} again
\end{itemize}

\begin{tcolorbox}[breakable,colback=teal!4,colframe=teal!35!black,title={Before you move on}]
What does \te{చురుకైన} mean? (\textquotedblleft{}Active, alert\textquotedblright{}.) Say it once more.
\end{tcolorbox}
\section[\texorpdfstring{\te{అల్లరి} (allari)}{అల్లరి (allari)}]{\te{అల్లరి} (allari) — naughty; mischief}

\label{lesson:TE-C297-allari}



\begin{quote}
Before the new one: say the Telugu for stingy, then the Telugu for active.
\end{quote}

\subsection*{You'll want to know: \te{అల్లరి}}
\textbf{\te{అల్లరి}} — \emph{allari} — \textquotedblleft{}naughty; mischief\textquotedblright{}.

\begin{grammarlens}[title={What you've built}]
One more word for this chapter.
\end{grammarlens}

\subsection*{Your turn}
\begin{itemize}
\raggedright
  \item \emph{Say it:} \emph{allari}
  \item \emph{Say it:} \emph{allari}, once more
  \item \emph{Say it:} say \emph{allari}
  \item \emph{Recall:} say the Telugu for stingy, then the Telugu for active, then say \emph{allari} again
\end{itemize}

\begin{tcolorbox}[breakable,colback=teal!4,colframe=teal!35!black,title={Before you move on}]
What does \te{అల్లరి} mean? (\textquotedblleft{}Naughty; mischief\textquotedblright{}.) Say it once more.
\end{tcolorbox}
\section[\texorpdfstring{\te{నిజాయితీ} (nijāyitī)}{నిజాయితీ (nijāyitī)}]{\te{నిజాయితీ} (nijāyitī) — honesty}

\label{lesson:TE-C297-nijayiti}



\begin{quote}
Before the new one: say the Telugu for active, then the Telugu for naughty.
\end{quote}

\subsection*{You'll want to know: \te{నిజాయితీ}}
\textbf{\te{నిజాయితీ}} — \emph{nijāyitī} — \textquotedblleft{}honesty\textquotedblright{}.

\begin{grammarlens}[title={What you've built}]
One more word for this chapter.
\end{grammarlens}

\subsection*{Your turn}
\begin{itemize}
\raggedright
  \item \emph{Say it:} \emph{nijāyitī}
  \item \emph{Say it:} \emph{nijāyitī}, once more
  \item \emph{Say it:} say \emph{nijāyitī}
  \item \emph{Recall:} say the Telugu for active, then the Telugu for naughty, then say \emph{nijāyitī} again
\end{itemize}

\begin{tcolorbox}[breakable,colback=teal!4,colframe=teal!35!black,title={Before you move on}]
What does \te{నిజాయితీ} mean? (\textquotedblleft{}Honesty\textquotedblright{}.) Say it once more.
\end{tcolorbox}
